% Generated draft: requires mathematical and source review.
\documentclass[11pt]{article}
\usepackage[margin=25mm]{geometry}
\usepackage{fontspec}
\setmainfont{Latin Modern Roman}
\usepackage{amsmath,amssymb,mathtools,unicode-math}
\setmathfont{Latin Modern Math}
\usepackage{enumitem,needspace,xurl,hyperref}
\setlength{\parindent}{0pt}
\setlength{\parskip}{5pt}
\setlength{\emergencystretch}{3em}
\newcommand{\N}{\mathbb{N}}
\newcommand{\Z}{\mathbb{Z}}
\newcommand{\Q}{\mathbb{Q}}
\newcommand{\R}{\mathbb{R}}
\newcommand{\C}{\mathbb{C}}
\newcommand{\F}{\mathbb{F}}
\newcommand{\A}{\mathbb{A}}
\newcommand{\PP}{\mathbb{P}}
\newcommand{\E}{\mathbb{E}}
\newcommand{\eps}{\varepsilon}
\newcommand{\dd}{\,\mathrm{d}}
\begin{document}
\section*{UnsolvedMath ID: 30004554}
\textbf{Source problem code:} catalogue entry
\subsection{Definitions and mathematical statement}
A Tauberian theorem derives a conclusion about an underlying mathematical object from information about a transformed or averaged version of that object, together with additional hypotheses. A remainder theorem supplies an explicit quantitative error estimate associated with such a conclusion. Here, proof mining refers to proof-theoretic analysis aimed at extracting quantitative information from proofs; the Dialectica interpretation is the example named in the supplied record.

The supplied research question is:
\begin{quote}
Can proof-theoretic methods such as the Dialectica interpretation produce new remainder theorems in Tauberian theory, and are there abstract quantitative metatheorems that unify and generalize related Tauberian theorems?
\end{quote}

This asks for both new quantitative Tauberian results obtained by proof-theoretic methods and abstract metatheorems explaining quantitative conclusions across related Tauberian results. The record does not specify the underlying objects or their domains, transforms or averaging operations, Tauberian hypotheses, error quantities, admissible parameter dependence, or intended family of related theorems. It also does not specify a formal theory, a class of proofs, or the logical form and quantitative conclusions of the proposed metatheorems. Consequently, the question is a research programme rather than a fully specified quantified proposition. No particular choice of these missing data is imposed here, and no existence theorem is asserted.
\subsection{Short English statement}
The question asks whether proof mining, including the Dialectica interpretation, can yield new explicit Tauberian error bounds and general quantitative metatheorems encompassing related results. The supplied record leaves the relevant theorem classes, hypotheses, and quantitative targets unspecified, so a unique formal proposition cannot yet be stated.
\subsection{Sources}
\begin{itemize}
\item Local registry: \texttt{data/unsolvedmath/problems.json}, permanent UnsolvedMath ID \texttt{30004554}, title ``Proof Mining for Quantitative Tauberian Remainders,'' source problem code \texttt{catalogue entry}. The supplied record belongs to the Oberwolfach Reports Open Problems collection and gives proposed year 2020. No individual proposer is supplied.
\item Supplied report citation: \emph{Mathematical Logic: Proof Theory, Constructive Mathematics} (2021), Oberwolfach Reports, pp.\ 1693--1757, DOI \texttt{10.4171/owr/2020/34}. \url{https://doi.org/10.4171/owr/2020/34}. The background identifies Oberwolfach Reports as documentation of mathematical workshops and their open-problem discussions; the supplied literature-source label is ``Oberwolfach Report 34/2020.'' No passage locating this particular question within the report is supplied.
\item Supplied literature reference: ``Powell, 2023.'' \url{https://doi.org/10.1016/j.apal.2022.103231}. No article title, theorem number, or detailed account of its relationship to the question is supplied.
\item Supplied literature assessment, checked \texttt{2026-08-22}: the record labels the problem ``partially solved'' and states that proof mining has produced concrete Tauberian remainder theorems and rates for specific results, while the requested general remainder theory and unifying quantitative metatheorems remain open. This is supplied dated curation derived from the OWR material, not an independently verified literature assessment. The background explicitly requests independent verification of time-sensitive status and cited claims.
\end{itemize}
\end{document}
